\section[APP levantada y terna espectral en dimensión 11]{Levantamiento operatorial de APP y terna espectral de la interfaz supersimétrica en dimensión \texorpdfstring{\(11\)}{11}}
\label{sec:sucesor83-terna-espectral-teoria-m}

Las pantallas \(11\to10\), la reducción isotrópica \(26\to24\) y la
involución entre momento y enrollamiento proporcionan el soporte algebraico
exacto de la interfaz. La dinámica física añade un espacio de estados, un
Hamiltoniano y una supercarga. Estos datos se reúnen en la terna
\[
 \mathfrak M
 =\bigl(\mathcal H_{\mathrm{phys}},H,Q\bigr),
\]
donde \(\mathcal H_{\mathrm{phys}}\) es un espacio de Hilbert
\(\mathbb Z_2\)-graduado,
\[
 \mathcal H_{\mathrm{phys}}
 =\mathcal H_+\oplus\mathcal H_-,
\]
\(H\) es un operador autoadjunto no negativo compatible con los transportes
y \(Q\) es un operador impar. La relación de supercarga es
\begin{equation}
 Q^2=H,
 \qquad
 [H,Q]=0.
 \label{eq:sucesor83-supercarga-cuadrado-hamiltoniano}
\end{equation}
La igualdad implica que \(Q\) intercambia los sectores graduados y que el
Hamiltoniano preserva cada uno de ellos.

\subsection{Levantamiento operatorial del estado APP--TRIT--TPK}

La APP enriquecida conserva simultáneamente los residuos y cocientes de las
hojas aditiva y multiplicativa:
\[
 \mathcal A_9^\sharp
 =\bigl(R^+,Q^+;R^\times,Q^\times\bigr).
\]
TRIT añade orientación temporal y el Topological Phase Kernel conserva
fase, ruta, acarreo, frontera y memoria. Sea
\(\mathcal H_{\mathrm{HMT}}\) el espacio lineal generado por estos estados
enriquecidos. Una realización operatorial compatible es un mapa
\[
 W:\mathcal H_{\mathrm{HMT}}
 \longrightarrow\mathcal H_{\mathrm{phys}}
\]
que preserva la graduación, transporta los observables de residuo y cociente
y entrelaza la involución bidireccional con la dualidad de radio. Para los
operadores definidos en ambos dominios, la compatibilidad exige
\[
 W\,U_T^{\mathrm{HMT}}
 =U_T^{\mathrm{phys}}\,W,
 \qquad
 W\,H_{\mathrm{HMT}}
 =H\,W.
\]
Esta flecha no identifica por decreto los dos espacios: especifica el dato
que debe preservar la genealogía al pasar de la construcción discreta a su
realización física.

\subsection{Compactificación y codominio físico undecadimensional}

La realización en dimensión \(11\) tiene por codominio un espacio
\[
 \mathcal X_{11}^{\mathrm{phys}}
\]
dotado de métrica \(g_{MN}\), \(3\)-forma \(C_{MNP}\), gravitino
\(\psi_M\), acción y transformaciones de supersimetría. La costura
\(\ell_K\) produce la descomposición exacta
\[
 V_{11}=\ell_K\oplus V_{10},
\]
y la compactificación debe inducir sobre
\(\mathcal H_{\mathrm{phys}}\) la descomposición espectral correspondiente.
La dualidad \(T\) exige, además,
\[
 U_T H(R)U_T^{-1}=H(R^{-1}).
\]
La superficie de mundo de la cuerda y los volúmenes de mundo de las branas
aportan entonces las realizaciones extendidas del mismo antecedente
genealógico.

\subsection{Compatibilidad espectral de \((\mathcal H_{\mathrm{phys}},H,\mathcal Q)\) con \(\mathcal Q^2=H\), dualidad T y compactificación \(11\to10\)}

Una realización satisface la interfaz cuando conserva conjuntamente:
\begin{enumerate}
\item la gradación de \(\mathcal H_{\mathrm{phys}}\) y la identidad
      \(Q^2=H\);
\item la acción de Polyakov y las condiciones de frontera de la superficie
      de mundo;
\item la polarización \(11\to10\) y la reducción isotrópica \(26\to24\);
\item el intercambio de momento y enrollamiento bajo \(R\leftrightarrow
      R^{-1}\);
\item el descenso de los campos \(g_{MN}\), \(C_{MNP}\) y \(\psi_M\) al
      límite compactificado;
\item la memoria de fase, ruta, orientación, acarreo y frontera del estado
      APP--TRIT--TPK.
\end{enumerate}
Las identidades de pantallas y la involución de radio son exactas en sus
dominios algebraicos. La terna
\((\mathcal H_{\mathrm{phys}},H,Q)\), la acción y los campos
undecadimensionales fijan, de forma separada y explícita, los datos de la
interfaz física.
